\section{Related Work}
\label{sec:related}

\paragraph{AIGC Detection.}
Supervised AIGC detectors, typically BERT-based classifiers trained
on human/machine text corpora~\citep{devlin2019bert, he2023mgtbench},
are complemented by zero-shot methods exploiting probability
curvature~\citep{mitchell2023detectgpt} or
watermarking~\citep{kirchenbauer2023watermark}.
However, \citet{sadasivan2023aigenerated} show that simple
paraphrasing dramatically reduces detection accuracy, and
cross-domain transfer remains
fragile~\citep{he2023mgtbench, yang2024survey}.

\paragraph{Text Style Transfer and Adversarial Probing.}
Style transfer methods---variational autoencoders, back-translation,
retrieve-and-edit pipelines, and LLM
rewriting~\citep{lyu2021style}---all operate on discrete tokens and
cannot continuously modulate transformation intensity.
Token-level adversarial attacks such as
TextFooler~\citep{jin2020bert} flip classifier predictions through
lexical substitution but likewise lack continuous controllability.
\citet{krishna2024paraphrasing} show that paraphrasing can evade
detectors, though retrieval-based defenses offer partial mitigation.
All prior methods thus lack a continuous diagnostic axis over
transformation intensity---a capability that continuous
embedding-space generation uniquely affords.

\paragraph{Diffusion and Flow Matching for Text.}
\citet{li2022diffusion} propose Diffusion-LM for controllable
generation in embedding space; \citet{lou2024discrete} develop MDLM
for discrete diffusion. \citet{chen2026langflow} show that continuous
flow matching rivals discrete diffusion in language modeling.
Existing methods address unconditional or class-conditional
generation; source-conditioned style transfer remains
unexplored---a gap \styleshield{} fills.